\midsectiontitle{Combate à Distância}

Em \textit{Overgrown}, atacar à distância não é o caminho do covarde, é a arte de quem compreende que sobreviver é tão valioso quanto matar. Arqueiros, besteiros e artilheiros dominam o espaço entre eles e o perigo. As \textbf{Habilidades de Distância} são técnicas ativas que consomem \textbf{Pontos de Combate (PC)} para transformar cada projétil em uma decisão tática precisa.

\begin{rpgboxtips}{Alcance e Cobertura}
Ataques à distância sofrem penalidades de cobertura parcial e total. Certas habilidades desta lista existem especificamente para ignorar essas penalidades, aprenda quando utilizá-las ou o inimigo que se abriga atrás de uma rocha será seu maior problema.
\end{rpgboxtips}

\section*{Habilidades}

\combatSkill{Mira Focada}{2}{Gaste sua \textbf{Ação de Movimento} para \textit{travar} o alvo em sua linha de mira. Seu próximo ataque contra ele \textbf{ignora cobertura parcial}. A preparação deliberada transforma incerteza em certeza.}

\combatSkill{Flecha de Arrasto}{2}{Ao acertar um inimigo que já está \link{status:sangramento}{Sangrando}, a flecha possui farpas internas que dilaceram o ferimento na saída. O alvo sofre \textbf{o dano do sangramento imediatamente}, além do dano normal do ataque.}

\combatSkill{Tiro de Supressão}{3}{Use sua \textbf{Reação} para disparar um projétil no momento exato em que um alvo começa a conjurar. Ao acertar, a conjuração é \textbf{interrompida} e a ação gasta é perdida.}

\combatSkill{Ponto Cego}{2}{Ao atacar um alvo que está \textbf{engajado em combate corpo a corpo com um aliado}, você gasta PC para encontrar uma brecha. O ataque \textbf{ignora a Esquiva} do alvo (apenas a Resistência é aplicada) e \textbf{elimina o risco de acertar seu aliado}.}

\combatSkill{Tiro de Reação}{3}{Quando um inimigo entrar no seu \textbf{alcance de tiro curto}, use sua \textbf{Reação} para disparar imediatamente. Se acertar, o inimigo \textbf{interrompe o movimento} e permanece onde está pelo restante do turno.}

\combatSkill{Sniper: Olho por Olho}{4}{Após receber dano de um \textbf{ataque à distância}, seu próximo ataque contra quem te atirou ganha \textbf{+4 na Margem de Ameaça}. A vingança aguça a mira.}

\combatSkill{Sniper: Olhos de Águia}{X (máx. 10)}{Para cada PC gasto, você aumenta o \textbf{alcance efetivo} do seu ataque em \textbf{+1 metro}, ignorando penalidades por distância acima do alcance base. Custo máximo de 10 PC por ataque.}

\combatSkill{Tiro de Aviso}{1}{Você atira propositalmente perto da cabeça do alvo. Realize um teste de \textbf{Intimidação usando Graça} contra a Temperança do alvo. Em caso de sucesso, o alvo hesita ou recua.}

\combatSkill{Recarregar Sob Pressão}{3}{Ignore a \textbf{penalidade de tempo de recarga} de armas pesadas como bestas pesadas ou armas de fogo lentas, podendo atirar \textbf{no mesmo turno} em que recarregou.}

\combatSkill{Mira de Precisão}{X}{Para cada PC gasto, você acumula \textbf{+1 de acerto} para um único tiro neste turno. Concentração total convertida em precisão absoluta.}

\combatSkill{Olhar de Rapina}{4 por rodada}{Enquanto estiver \textbf{parado}, gaste PC por rodada para focar o mesmo alvo. A cada rodada consecutiva, sua \textbf{Margem de Ameaça} aumenta em \textbf{+2} cumulativamente. Se você se mover ou trocar de alvo, o bônus acumulado zera.}

\combatSkill{Tiro Ricochete}{4}{O projétil \textbf{ricocha uma vez} em uma superfície antes de atingir o alvo, contornando cobertura parcial ou total. Ricochetes não podem resultar em crítico.
\par\textit{\footnotesize \textbf{Aprimorado:} O tiro ricocheteado \textbf{pode resultar em crítico} normalmente.}}

\combatSkill{Tiro Destroçante}{3}{Gaste PC para converter todo o \textbf{dano do ataque} em dano direto à \textbf{Armadura do alvo} ao invés de sua vida, reduzindo os PV da armadura e tornando-a menos efetiva para os demais combatentes.}

\combatSkill{Desengajar}{2}{Você chuta o inimigo mais próximo para se \textbf{propulsionar para trás}, ganhando distância sem provocar \textbf{ataques de oportunidade}. Ideal para manter o alcance seguro.}

\combatSkill{Lobo Solitário}{5}{Você fecha sua mente para o grupo. Curas e Buffs aliados são \textbf{reduzidos à metade} em você, mas todos os seus \textbf{ataques e habilidades} ganham \textbf{+1D6 de dano} adicional.
\par\textit{\footnotesize \textbf{Aprimorado:} O bônus de dano aumenta para \textbf{+3D6} em todos os ataques e habilidades.}}